\begin{table}[t]
\centering
\caption{Table A — i\_rehearsal\_v2（接入/harness 验证，不是实验数据）　脚注（N-348，2026-09-10 用户裁定）：适配赛道 v1.0-adapt 的题源是**出集规定题的 oracle 产物**（探针题不入），这些产物随适配 bundle 进入执行面。**跑过适配赛道的被测方，主赛道这些题算「可能已见过答案」**；口径与逐题清单见 ops/specs/fairness\_protocol.md §7 与 ops/reports/known\_limits\_v1.md。}
\label{tab:a-i_rehearsal_v2}
\begin{tabular}{llrrrrrrrrr}
\toprule
config\_id & arm & SR & pass@1 & pass\^{}3 & ProgressRate & Steps & \$ & Latency & Recov & 越权率 \\
\midrule
\multicolumn{11}{l}{\textit{参照臂（baseline）}} \\
cfg-codex-deepseek & open & 0.667 & 0.000 & --- & 0.667 & 73.000 & 1.074 & 906.705 & --- & --- \\
\midrule
\multicolumn{11}{l}{\textit{协议臂（protocol）—— 与参照臂的差异是投放的工件}} \\
cfg-codex-deepseek & strict & 0.333 & 0.000 & --- & 0.333 & 93.667 & 2.067 & 1504.097 & --- & --- \\
\bottomrule
\end{tabular}
\end{table}

% set_version = 1.0.15
% reference_version = r1.0.22
% protocol_version = geneprotocol_v1@d6fbcaa08302（裸臂在协议轴上记 `geneprotocol_v1@none` —— 它们不投放协议工件，与协议臂并排不算混轴）
% channel = public
% 记录数 6；筛选 batch=i_rehearsal_v2
